\section{AIME I 1993}
        \begin{enumerate}
        	\item (\textit{AIME I 1993}) How many even integers between 4000 and 7000 have four different digits? 


 
 \textbf{Answer: }728



\textbf{Solution 1}

The thousands digit is $\in \{4,5,6\}$. 

 Case $1$: Thousands digit is even

 $4, 6$, two possibilities, then there are only $\frac{10}{2} - 1 = 4$ possibilities for the units digit. This leaves $8$ possible digits for the hundreds and $7$ for the tens places, yielding a total of $2 \cdot 8 \cdot 7 \cdot 4 = 448$.

 
Case $2$: Thousands digit is odd

 $5$, one possibility, then there are $5$ choices for the units digit, with $8$ digits for the hundreds and $7$ for the tens place. This gives $1 \cdot 8 \cdot 7 \cdot 5= 280$ possibilities.

 Together, the solution is $448 + 280 = \boxed{728}$.

 



\textbf{Solution 2 by PEKKA(more elaborate)}

Firstly, we notice that the thousands digit could be $4$, $5$ or $6$.Since the parity of the possibilities are different, we cannot cover all cases in one operation. Therefore, we must do casework.Case $1$: 

 Here we let thousands digit be $4$.

 4 \_ \_ \_ 

 We take care of restrictions first, and realize that there are 4 choices for the last digit, namely $2$,$6$,$8$ and $0$. Now that there are no restrictions we proceed to find that there are $8$ choices for the hundreds digit and $7$ choices for the tens digit. Therefore we have $8 \cdot 7 \cdot 4 = 224$ numbers that satisfy the conditions posed by the problem.

 Case $2$

 Here we let thousands digit be $5$.

 5 \_ \_ \_ 

 Again, we take care of restrictions first. This time there are 5 choices for the last digit, which are all the even numbers because 5 is odd.

 Now that there are no restrictions we proceed to find that there are $8$ choices for the hundreds digit and $7$ choices for the tens digit. Therefore we have $8 \cdot 7 \cdot 5 = 280$ numbers that satisfy the conditions posed by the problem.

 Case $3$

 Here we let thousands digit be $6$.

 6 \_ \_ \_ 

 Once again, we take care of restrictions first. Since 6 is even,  there are 4 choices for the last digit, namely $2$,$6$,$8$ and $0$

 Now we proceed to find that there are $8$ choices for the hundreds digit and $7$ choices for the tens digit. Therefore we have $8 \cdot 7 \cdot 4 = 224$ numbers that satisfy the conditions posed by the problem.

 Now that we have our answers for each possible thousands place digit, we add up our answers and get $224+280+224$= $\boxed{728}$~PEKKA

 


	\item (\textit{AIME I 1993}) During a recent campaign for office, a candidate made a tour of a country which we assume lies in a plane. On the first day of the tour he went east, on the second day he went north, on the third day west, on the fourth day south, on the fifth day east, etc. If the candidate went $\frac{n^{2}}{2}$ miles on the $n^{\mbox{th}}_{}$ day of this tour, how many miles was he from his starting point at the end of the $40^{\mbox{th}}_{}$ day? 


 
 \textbf{Answer: }580



\textbf{Solution}

On the first day, the candidate moves $[4(0) + 1]^2/2\ \text{east},\, [4(0) + 2]^2/2\ \text{north},\, [4(0) + 3]^2/2\ \text{west},\, [4(0) + 4]^2/2\ \text{south}$, and so on. The east/west displacement is thus $\frac{1^2 - 3^2 + 5^2 \ldots +37^2 - 39^2}{2} = \left|\sum_{i=0}^9 \frac{(4i+1)^2}{2} - \sum_{i=0}^9 \frac{(4i+3)^2}{2}\right|$. Applying \href{https://artofproblemsolving.com/wiki/index.php?title=Difference\_of\_squares}{difference of squares}, we see that 
 \begin{align*} \left|\sum_{i=0}^9 \frac{(4i+1)^2 - (4i+3)^2}{2}\right| &= \left|\sum_{i=0}^9 \frac{(4i+1+4i+3)(4i+1-(4i+3))}{2}\right|\\ &= \left|\sum_{i=0}^9 -(8i+4) \right|. \end{align*}
Similarly, the north/south displacement is 
 \[\left|\sum_{i=0}^9 \frac{(4i+2)^2}{2} - \sum_{i=0}^9 \frac{(4i+4)^2}{2}\right| = \left|\sum_{i=0}^9 -(8i+6) \right|.\]Since $\sum_{i=0}^{9} i = \frac{9(10)}{2} = 45$, the two distances evaluate to $8(45) + 10\cdot 4 = 400$ and $8(45) + 10\cdot 6 = 420$. By the \href{https://artofproblemsolving.com/wiki/index.php?title=Pythagorean\_Theorem}{Pythagorean Theorem}, the answer is $\sqrt{400^2 + 420^2} = 29 \cdot 20 = \boxed{580}$.

 


	\item (\textit{AIME I 1993}) The table below displays some of the results of last summer's Frostbite Falls Fishing Festival, showing how many contestants caught $n\,$ fish for various values of $n\,$. 


 $\begin{array}{|c|c|c|c|c|c|c|c|c|} \hline n & 0 & 1 & 2 & 3 & \dots & 13 & 14 & 15 \\ \hline \text{number of contestants who caught} \ n \ \text{fish} & 9 & 5 & 7 & 23 & \dots & 5 & 2 & 1 \\ \hline \end{array}$

 In the newspaper story covering the event, it was reported that 


 
            \begin{enumerate}[label=(\roman*)]
                \item  the winner caught $15$ fish;
\item  those who caught $3$ or more fish averaged $6$ fish each;
\item  those who caught $12$ or fewer fish averaged $5$ fish each.
            \end{enumerate}
            What was the total number of fish caught during the festival?


 
 \textbf{Answer: }943



\textbf{Solution 1}

Suppose that the number of fish is $x$ and the number of contestants is $y$. The $y-(9+5+7)=y-21$ fishers that caught $3$ or more fish caught a total of $x - \left(0\cdot(9) + 1\cdot(5) + 2\cdot(7)\right) = x - 19$ fish. Since they averaged $6$ fish, 

 $6 = \frac{x - 19}{y - 21} \Longrightarrow x - 19 = 6y - 126.$

  Similarily, those who caught $12$ or fewer fish averaged $5$ fish per person, so 

 $5 = \frac{x - (13(5) + 14(2) + 15(1))}{y - 8} = \frac{x - 108}{y - 8} \Longrightarrow x - 108 = 5y - 40.$

  Solving the two equation system, we find that $y = 175$ and $x = \boxed{943}$, the answer.

 



\textbf{Solution 2}

Let $f$ be the total number of fish caught by the contestants who didn't catch $0, 1, 2, 3, 13, 14$, or $15$ fish and let $a$ be the number of contestants who didn't catch $0, 1, 2, 3, 13, 14$, or $15$ fish. From $\text{(b)}$, we know that $\frac{69+f+65+28+15}{a+31}=6\implies f=6a+9$. From $\text{(c)}$ we have $\frac{f+69+14+5}{a+44}=5\implies f=5a+132$. Using these two equations gets us $a=123$. Plug this back into the equation to get $f=747$. Thus, the total number of fish caught is $5+14+69+f+65+28+15=\boxed{943}$ - Heavytoothpaste

 



\textbf{Solution 3}

Let $x$ be the average number of fish caught among $n$ of those who caught $4$ to $12$ fishes. According to statement B:
 \[\frac{xn + 3 \cdot 23 + 13 \cdot 5 + 14 \cdot 2 + 15}{n + 31} = 6\]
 \[\implies xn - 6n = 9\]According to statement C:
 \[\frac{xn + 0 \cdot 9 + 1 \cdot 5 + 2 \cdot 7 + 3 \cdot 23}{n + 44} = 5\]
 \[\implies xn - 5n = 132\]We then solve for $n$ and then $xn$. Note $xn$ is the number of fish caught by the $n$ people we defined above.
 \[n = 123\]
 \[xn = 132 + 5(123) = 747\]Thus, the total number of fish caught is
 \[747 + 1 \cdot 5 + 2 \cdot 7 + 3 \cdot 23 + 13 \cdot 5 + 14 \cdot 2 + 15 \cdot 1 = \boxed{943}\] ~Totient4Breakfast

 


	\item (\textit{AIME I 1993}) How many ordered four-tuples of integers $(a,b,c,d)\,$ with $0 < a < b < c < d < 500\,$ satisfy $a + d = b + c\,$ and $bc - ad = 93\,$?


 
 \textbf{Answer: }870



\textbf{Solution 1}

Let $k = a + d = b + c$ so $d = k-a, b=k-c$. It follows that $(k-c)c - a(k-a) = (a-c)(a+c-k) = (c-a)(d-c) = 93$. Hence $(c - a,d - c) = (1,93),(3,31),(31,3),(93,1)$.

 Solve them in terms of $c$ to get $(a,b,c,d) = (c - 93,c - 92,c,c + 1),$ $(c - 31,c - 28,c,c + 3),$ $(c - 1,c + 92,c,c + 93),$ $(c - 3,c + 28,c,c + 31)$. The last two solutions don't follow $a < b < c < d$, so we only need to consider the first two solutions.

 The first solution gives us $c - 93\geq 1$ and $c + 1\leq 499$ $\implies 94\leq c\leq 498$, and the second one gives us $32\leq c\leq 496$.

 So the total number of such quadruples is $405 + 465 = \boxed{870}$.

 



\textbf{Solution 2}

Let $b = a + m$ and $c = a + m + n$. From $a + d = b + c$, $d = b + c - a = a + 2m + n$. 

 Substituting $b = a + m$, $c = a + m + n$, and $d = b + c - a = a + 2m + n$ into $bc - ad = 93$,
 \[bc - ad = (a + m)(a + m + n) - a(a + 2m + n) = m(m + n). = 93 = 3(31)\]Hence, $(m,n) = (1,92)$ or $(3,28)$. 

 For $(m,n) = (1,92)$, we know that $0 < a < a + 1 < a + 93 < a + 94 < 500$, so there are $405$ four-tuples. For $(m,n) = (3,28)$, $0 < a < a + 3 < a + 31 < a + 34 < 500$, and there are $465$ four-tuples. In total, we have $405 + 465 = \boxed{870}$ four-tuples.

 



\textbf{Solution 3}

Square both sides of the first equation in order to get $bc$ and $ad$ terms, which we can plug $93$ in for.
 \begin{align*} (a+d)^2 = (b+c)^2 &\implies a^2 + 2ad + d^2 = b^2 + 2bc + c^2 \\ &\implies 2bc-2ad = a^2-b^2 + d^2-c^2 \\ &\implies 2(bc-ad) = (a-b)(a+b)+(d-c)(d+c) \end{align*}
We can plug $93$ in for $bc - ad$ to get $186$ on the left side, and also observe that $a-b = c-d$ after rearranging the first equation. Plug in $c-d$ for $a-b$.

 $186 = (c-d)(a+b) + (d-c)(d+c) \implies 186 = -(d-c)(a+b) + (d-c)(d+c) \implies 186 = (d-c)(d+c-a-b)$

 Now observe the possible factors of $186$, which are $1 \cdot 186, 2\cdot 93, 3 \cdot 62, 6\cdot 31$. $(d-c)$ and $(d+c-a-b)$ must be factors of $186$, and $(d+c-a-b)$ must be greater than $(d-c)$.

 $1 \cdot 186$ work, and yields $405$ possible solutions.$2 \cdot 93$ does not work, because if $c-d = 2$, then $a+b$ must differ by 2 as well, but an odd number $93$ can only result from two numbers of different parity. $c-d$ will be even, and $a+b$ will be even, so $c+d - (a+b)$ must be even.$3 \cdot 62$ works, and yields $465$ possible solutions, while $6 \cdot 31$ fails for the same reasoning above.

 Thus, the answer is $405 + 465 = \boxed{870}$




\textbf{Solution 4}

Add the two conditions together to get $a+d+ad+93=b+c+bc$. Rearranging and factorising with SFFT, $(a+1)(d+1)+93=(b+1)(c+1)$. This implies that for every quadruple $(a,b,c,d)$, we can replace $a\longrightarrow a+1$, $b\longrightarrow b+1$, etc. and this will still produce a valid quadruple. This means, that we can fix $a=1$, and then just repeatedly add $1$ to get the other quadruples. 

 Now, our conditions are $b+c=d+1$ and $bc=d+93$. Replacing $d$ in the first equation, we get $bc-b-c=92$. Factorising again with SFFT gives $(b-1)(c-1)=93$. Since $b<c$, we have two possible cases to consider.

 Case 1: $b=2$, $c=94$. This produces the quadruple $(1,2,94,95)$, which indeed works.

 Case 2: $b=4$, $c=32$. This produces the quadruple $(1,4,32,35)$, which indeed works. 

 Now, for case 1, we can add $1$ to each term exactly $404$ times (until we get the quadruple $(405,406,498,499)$), until we violate $d<500$. This gives $405$ quadruples for case 1.

 For case 2, we can add $1$ to each term exactly $464$ times (until we get the quadruple $(465,468,496,499)$). this gives $465$ quadruples for case 2.

 In conclusion, having exhausted all cases, we can finish. There are hence $405+465=\boxed{870}$ possible quadruples.

 



\textbf{Solution 5}

Let $r = d-c$. From the equation $a+d = b+c$, we have 
 \[r = d-c = b-a ,\] so $b = a+r$ and $c = d-r$. We then have 
 \[93 = (a+r)(d-r) - ad = rd - ra - r^2 = r(d-a-r) .\] Since $c > b$, $d-r > a+r$, or $d-a-r > r$. Since the prime factorization of 93 is $3 \cdot 31$, we must either have $r=1$ and $d-a-r = 93$, or $r=3$ and $d-a-r = 31$. We consider these cases separately.

 If $r=1$, then $d-a = 94$, $b= a+1$, and $c = d-1$. Thus $d$ can be any integer between 95 and 499, inclusive, and our choice of $d$ determines the four-tuple $(a,b,c,d)$. We therefore have $499-95+1 = 405$ possibilities in this case.

 If $r=3$, then $d-a = 34$, $b = a+3$, and $c=d-3$. Thus $d$ can be any integer between 35 and 499, inclusive, and our choice of $d$ determines the four-tuple $(a,b,c,d)$, as before. We therefore have $499-35+1 = 465$ possibilities in this case.

 Since there are 405 possibilities in the first case and 465 possibilities in the second case, in total there are $405 + 465 = \boxed{870}$ four-tuples.

 



\textbf{Solution 6}

Assume $d = x+m, a = x-m, c = x+n$, and $b = x-n$. This clearly satisfies the condition that $a+d = b+c$ since ($2x = 2x$) . Now plug this into $bc-ad = 93$. You get $(x+n)(x-n) - (x+m)(x-m) = 93 \Rightarrow m^2 - n^2 = 93 \Rightarrow (m-n)(m+n) = 93$

 Since $m>n$ (as given by the condition that $a<b<c<d$), $m+n>m-n$ and $m$ and $n$ are integers, there are two cases we have to consider since $93 = 3\cdot 31$. We first have to consider $m-n = 1, m+n = 93$, and then consider $m-n=3, m+n = 31$. 

 In the first case, we get $m=47, n=46$ and in the second case we get $m=17, n=14$. Now plug these values (in separate cases) back into $a,b,c,d$. Since the only restriction is that all numbers have to be greater than $0$ or less than $500$, we can write two inequalities. Either $x+47 < 500, x-47 > 0$, or $x+17 < 500, x-17 > 0$ (using the inequalities given by $d$ and $a$, and since $b$ and $c$ are squeezed in between $d$ and $a$, we only have to consider these two inequalities).

 This gives us either $47 < x < 453$ or $17 < x < 483$, and using simple counting, there are $405$ values for $x$ in the first case and $465$ values for $x$ in the second case, and hence our answer is $405+465 = \boxed{870}$



	\item (\textit{AIME I 1993}) Let $P_0(x) = x^3 + 313x^2 - 77x - 8\,$.  For integers $n \ge 1\,$, define $P_n(x) = P_{n - 1}(x - n)\,$.  What is the coefficient of $x\,$ in $P_{20}(x)\,$?


 
 \textbf{Answer: }763



\textbf{Solution 1}

Notice that 
 \begin{align*}P_{20}(x) &= P_{19}(x - 20)\\ &= P_{18}((x - 20) - 19)\\ &= P_{17}(((x - 20) - 19) - 18)\\ &= \cdots\\ &= P_0(x - (20 + 19 + 18 + \ldots + 2 + 1)).\end{align*}


 
Using the formula for the sum of the first $n$ numbers, $1 + 2 + \cdots + 20 = \frac{20(20+1)}{2} = 210$. Therefore, 
 \[P_{20}(x) = P_0(x - 210).\]

 Substituting $x - 210$ into the function definition, we get $P_0(x-210) = (x - 210)^3 + 313(x - 210)^2 - 77(x - 210) - 8$. We only need the coefficients of the linear terms, which we can find by the \href{https://artofproblemsolving.com/wiki/index.php?title=Binomial\_theorem}{binomial theorem}. 

 
            \begin{itemize}
        \item $(x-210)^3$ will have a linear term of ${3\choose1}210^2x = 630 \cdot 210x$.
        \item $313(x-210)^2$ will have a linear term of $-313 \cdot {2\choose1}210x = -626 \cdot 210x$.
        \item $-77(x-210)$ will have a linear term of $-77x$.
            \end{itemize}
            Adding up the coefficients, we get $630 \cdot 210 - 626 \cdot 210 - 77 = \boxed{763}$.

 



\textbf{Solution 2 (Overkill)}

Denote the coefficient of term $x^m$ of polynomial $P_n(x)$ as $c_{m, n}$. We can see that $c_{2, 0} = 313$, $c_{1, 0} = -77$, and $c_{0, 0} = -8$. Note that $(x-n)^3 = x^3-3nx^2+3n^2x-n^3$ and $(x-n)^2 = x^2-2nx+n^2$When substituting $x-1$ for $x$ in $P_0(x)$, we get that $P_1(x) = P_0(x-1) = (x^3-3x^2+3x-1) + 313(x^2-2x+1) -77(x-1) -8$. Observe that
 \[c_{2,1} = c_{2,0} - 3(1) = 313 - 3 = 310\]It is evident that
 \[c_{2,n} = c_{2,n-1} - 3n\]Define the generating function $C_2(X)$ as
 \[C_2(X) := \sum_{n \geq 0} c_{2,n}X^n\]By multiplying both sides of the previous recurrence relation and taking the sum of the terms $\forall n\geq 1$, we get that
 \[\implies \sum_{n\geq1} c_{2,n}X^n = \sum_{n\geq1}c_{2,n-1}X^n + 3\sum_{n\geq1}nX^n\]
 \[\implies \sum_{n\geq0} c_{2,n}X^n - c_{2,0} = X\sum_{n\geq0}c_{2,n}X^n + 3\sum_{n\geq1}nX^n\]
 \[\implies C_2(X) - 313 = XC_2(X) + 3\cdot \frac{X}{(1-X)^2}^*\]
 \[\implies C_2(X) = \frac{313}{1-X} - \frac{3X}{(1-X)^3}\]
 \[= \frac{313X^2 - 629X + 313}{(1-X)^3} = (313X^2 - 629X + 313) \cdot \sum_{n\geq0}{n+2 \choose 2}X^n\ ^\dagger\]
 \[\implies c_{2,n} = 313{n+2 \choose 2} - 629{n+1 \choose 2} + 313{n \choose 2}^\ddagger = \frac{-3n^2 - 3n + 626}{2}\]We can perform a similar analysis on $c_{1,n}$ to get the recurrence relation
 \[c_{1,n} = c_{1, n-1} + c_{2,n-1}\cdot(-2n) + 3n^2\]
 \[= c_{1, n-1} +3n^3 - 626n\] Define the generating function $C_1(X)$ as
 \[C_1(X) := \sum_{n\geq0}c_{1,n}X^n\]We can then perform this process again on our new recurrence relation:
 \[\implies \sum_{n\geq1} c_{1,n}X^n = \sum_{n\geq1}c_{1,n-1}X^n + 3\sum_{n\geq1}n^3X^n - 626\sum_{n\geq1}nX^n\]
 \[\implies \sum_{n\geq0} c_{1,n}X^n - c_{1,0} = X\sum_{n\geq0}c_{1,n}X^n + 3\sum_{n\geq1}n^3X^n - 626\sum_{n\geq1}nX^n\]
 \[\implies C_1(X) + 77 = XC_1(X) + 3\cdot \frac{X(X^2+4X+1)}{(1-X)^4} - 626\cdot \frac{X}{(1-X)^2}\]
 \[\implies C_1(X) = \frac{-77(1-X)^4 + 3X(X^2+4X+1) - 626X(1-X)^2}{(1-X)^5}\]
 \[= \frac{-77X^4-315X^3+802X^2-315X-77}{(1-X)^5}\]
 \[=(-77X^4-315X^3+802X^2-315X-77) \cdot \sum_{n\geq0}{n+5 \choose 5}X^n\]
 \[\implies c_{1,n} = -77{n+4 \choose 4} - 315{n+3 \choose 4} + 802{n+2 \choose 4} - 315{n+1 \choose 4} -77{n \choose 4}\]
 \[= \frac{3n^4+6n^3-1249n^2-1252n-308}{4}\]Finally, we can plug $n=20$ into our new explicit formula to get 
 \[c_{1,20} = \boxed{763}\]$^*$This can be calculated by differentiating the generating function of the sequence 1, 1, 1, ... $(\sum_{n\geq0}X^n)$ and multiplying by $X$.

 $^\dagger$This form can be found by applying Newton's generalized binomial theorem.

 $^\ddagger$This formula can be found by convolving the polynomial with the series.

 - MathCactus0\_0 (don't try this on a test unless you can't think of anything else!!!)

 


	\item (\textit{AIME I 1993}) What is the smallest positive integer than can be expressed as the sum of nine consecutive integers, the sum of ten consecutive integers, and the sum of eleven consecutive integers?


 
 \textbf{Answer: }495



\textbf{Solution 1}

Denote the first of each of the series of consecutive integers as $a,\ b,\ c$. Therefore, $n = a + (a + 1) \ldots (a + 8) = 9a + 36 = 10b + 45 = 11c + 55$. Simplifying, $9a = 10b + 9 = 11c + 19$. The relationship between $a,\ b$ suggests that $b$ is divisible by $9$. Also, $10b -10 = 10(b-1) = 11c$, so $b-1$ is divisible by $11$. We find that the least possible value of $b = 45$, so the answer is $10(45) + 45 = 495$.

 



\textbf{Solution 2}

Let the desired integer be $n$. From the information given, it can be determined that, for positive integers $a, \ b, \ c$:

 $n = 9a + 36 = 10b + 45 = 11c + 55$

 This can be rewritten as the following congruences:

 $n \equiv 0 \pmod{9}$
$n \equiv 5 \pmod{10}$
$n \equiv 0 \pmod{11}$

 Since 9 and 11 are relatively prime, n is a multiple of 99. It can then easily be determined that the smallest multiple of 99 with a units digit 5 (this can be interpreted from the 2nd congruence) is $\boxed{495}$




\textbf{Solution 3}

Let $n$ be the desired integer. From the given information, we have
 \begin{align*}9x &= a \\ 11y &= a \\ 10z + 5 &= a, \end{align*}
 here, $x,$ and $y$ are the middle terms of the sequence of 9 and 11 numbers, respectively. Similarly, we have $z$ as the 4th term of the sequence. Since, $a$ is a multiple of $9$ and $11,$ it is also a multiple of $\text{lcm}[9,11]=99.$ Hence, $a=99m,$ for some $m.$ So, we have $10z + 5 = 99m.$ It follows that $99(5) = \boxed{495}$ is the smallest integer that can be represented in such a way.

 



\textbf{Solution 4}

By the method in Solution 1, we find that the number $n$ can be written as $9a+36=10b+45=11c+55$ for some integers $a,b,c$. From this, we can see that $n$ must be divisible by 9, 5, and 11. This means $n$ must be divisible by 495. The only multiples of 495 that are small enough to be AIME answers are 495 and 990. From the second of the three expressions above, we can see that $n$ cannot be divisible by 10, so $n$ must equal $\boxed{495}$. Solution by Zeroman.

 



\textbf{Solution 5}

First note that the integer clearly must be divisible by $9$ and $11$ since we can use the "let the middle number be x" trick. Let the number be $99k$ for some integer $k.$ Now let the $10$ numbers be $x,x+1, \cdots x+9.$ We have $10x+45 = 99k.$ Taking mod $5$ yields $k \equiv 0 \pmod{5}.$ Since $k$ is positive, we take $k=5$ thus obtaining $99 \cdot 5 = \boxed{495}$ as our answer.

 


 



\textbf{Solution 6}

From the first equation solution 1, we get the three part equation 9x=10y+9=11z+19 after subtracting 36 from each of the equations 9x+36=10y+45=11z+55 knowing that the integer we are looking for has a last digit of 9 from the second part of our first equation, and also knowing that our mystery number must be 19 more the a multiple of 11 that ends in 0, the only way this could happen is if the number is able to be deduced in the form 110a+19, where a is between the range 1-9, inclusive. so, also knowing our mystery number is divisible by 9, our only answer being 129,239,349,459 ... etc. until 899. Brushing through our act of casework for this problem shows that the answer is 459... but we are not done yet! We have to add the 36 we subtracted at the beggining of our solution to get the equation 459+36, which is equal to our final answer... 495!!!

 \begin{verbatim}
	$michaellin16$
\end{verbatim}




	\item (\textit{AIME I 1993}) Three numbers, $a_1, a_2, a_3$, are drawn randomly and without replacement from the set $\{1, 2, 3,\ldots, 1000\}$. Three other numbers, $b_1, b_2, b_3$, are then drawn randomly and without replacement from the remaining set of $997$ numbers. Let $p$ be the probability that, after suitable rotation, a brick of dimensions $a_1 \times a_2 \times a_3$ can be enclosed in a box of dimension $b_1 \times b_2 \times b_3$, with the sides of the brick parallel to the sides of the box. If $p$ is written as a fraction in lowest terms, what is the sum of the numerator and denominator?


 
 \textbf{Answer: }005



\textbf{Solution 1 (Permutation)}

There is a total of $P(1000,6)$ possible ordered $6$-tuples $(a_1,a_2,a_3,b_1,b_2,b_3).$

 There are $C(1000,6)$ possible sets $\{a_1,a_2,a_3,b_1,b_2,b_3\}.$ We have five valid cases for the increasing order of these six elements:

 
            \begin{enumerate} 
        \item $aaabbb$

        \item $aababb$

        \item $aabbab$

        \item $abaabb$

        \item $ababab$

            \end{enumerate}
            Note that the $a$'s are different from each other, as there are $3!=6$ ways to permute them as $a_1,a_2,$ and $a_3.$ Similarly, the $b$'s are different from each other, as there are $3!=6$ ways to permute them as $b_1,b_2,$ and $b_3.$

 So, there is a total of $C(1000,6)\cdot5\cdot6^2$ valid ordered $6$-tuples. The requested probability is 
 \[p=\frac{C(1000,6)\cdot5\cdot6^2}{P(1000,6)}=\frac{C(1000,6)\cdot5\cdot6^2}{C(1000,6)\cdot6!}=\frac14,\] from which the answer is $1+4=\boxed{005}.$




\textbf{Solution 2 (Combination)}

Call the six numbers selected $x_1 > x_2 > x_3 > x_4 > x_5 > x_6$. Clearly, $x_1$ must be a dimension of the box, and $x_6$ must be a dimension of the brick. 

 
            \begin{itemize}
        \item If $x_2$ is a dimension of the box, then any of the other three remaining dimensions will work as a dimension of the box. That gives us $3$ possibilities.
        \item If $x_2$ is not a dimension of the box but $x_3$ is, then both remaining dimensions will work as a dimension of the box. That gives us $2$ possibilities.
        \item If $x_4$ is a dimension of the box but $x_2,x_3$ aren’t, there are no possibilities (same for $x_5$).
            \end{itemize}
            The total number of arrangements is ${6\choose3} = 20$; therefore, $p = \frac{3 + 2}{20} = \frac{1}{4}$, and the answer is $1 + 4 = \boxed{005}$.

 Note that the $1000$ in the problem, is not used, and is cleverly bypassed in the solution, because we can call our six numbers $x_1,x_2,x_3,x_4,x_5,x_6$ whether they may be $1,2,3,4,5,6$ in some order or $999,5,3,998,997,891$ in some order.

 



\textbf{Solution 3 (Hook Length Formula)}

Like Solution 2, call the six numbers selected $x_1 > x_2 > x_3 > x_4 > x_5 > x_6$. Using the Hook Length Formula, the number of valid configuration is $\frac{6!}{4\cdot3\cdot2\cdot3\cdot2}=5$. We proceed as Solution 2 does.

 



\textbf{Solution 4 (Catalan Numbers)}

As in Solutions 2 and 3, we let $x_1>x_2>x_3>x_4>x_5>x_6$ where each $x_i$ is a number selected.  It is clear that when choosing whether each number must be in the set with larger dimensions (the box) or the set with smaller dimensions (the brick) there must always be at least as many numbers in the former set as the latter.  We realize that this resembles Catalan numbers, where the indices of the numbers in the first set can be replaced with rising sections of a mountain, and the other indices representing falling sections of a mountain.  The formula for the $n$th Catalan number (where $n$ is the number of pairs of rising and falling sections) is 
 \[\frac{\binom{2n}{n}}{n+1}\]Thus, there are $\frac{\binom{6}{3}}{4}$ ways to pick which of $x_1,x_2,x_3,x_4,x_5,$ and $x_6$ are the dimensions of the box, and which are the dimensions of the brick, such that the condition is fulfilled.  There are $\binom{6}{3}$ total ways to choose which numbers make up the brick and box, so the probability of the condition being fulfilled is $\frac{\binom{6}{3}/4}{\binom{6}{3}}=\frac14\Longrightarrow \boxed{005}$.

 



\textbf{Solution 5 (1-1 Correspondence and Constructive Counting)}

In order (after rotation, of course), for the brick to fit in the box, we must have \( a_1 < b_1 \), \( a_2 < b_2 \), and \( a_3 < b_3 \). Also, to correct for the rotation, we must have that \( a_1 < a_2 < a_3 \), and \( b_1 < b_2 < b_3 \). Due to the fact that all \( a_1 \), \( a_2 \), \( a_3 \), \( b_1 \), \( b_2 \), and \( b_3 \) are distinct, and also the fact that each \( a_i \) and \( b_i \): \( 1 \le i \le 3 \) are a set of six distinct elements from the set \( \{1,2,3,\dots,100\} \), it suffices to just find the probability that \( a_1 < b_1 \), \( a_2 < b_2 \), and \( a_3 < b_3 \) from the set \( \{1,2,3,4,5,6\} \). If you aren't convinced, \( a_1 \), \( a_2 \), and \( a_3 \) are independently chosen from \( \{1,2,3,\dots,100\} \). Once chosen, \( b_1 \), \( b_2 \), and \( b_3 \) are chosen from the remaining 997 elements. Thus, any subset of six elements can satisfy the criteria of the dimensions of the brick and box. The simplest of these subsets is \( \{1,2,3,4,5,6\} \). Thus, there is a 1-1 correspondence (bijection).

 Then, we know that \( a_1 = 1 \) and \( b_3 = 6 \) (This is quite easy to see. Notice that \( a_1 \) is the smallest possible number and \( b_3 \) is the largest due to the simplicity after rotation of \( a_1 < b_1 \), \( a_2 < b_2 \), and \( a_3 < b_3 \), \( a_1 < a_2 < a_3 \), and \( b_1 < b_2 < b_3 \)). Thus, we must evaluate the possible cases for \( a_2 \), \( a_3 \), \( b_1 \), and \( b_2 \).

 Case 1: \( a_2 = 2 \). Then, there are three cases for \( a_3 \) (either 3, 4, or 5), in which \( b_1 \) and \( b_2 \) are uniquely determined. (if \( a_3 = 3 \), then \( b_1 = 4 \) and \( b_2 = 5 \); if \( a_3 = 4 \), \( b_1 = 3 \) and \( b_2 = 5 \); and if \( a_3 = 5 \), \( b_1 = 3 \) and \( b_2 = 4 \)).

 Case 2: \( a_2 = 3 \). Then, there are two cases for \( a_3 \) (4 or 5). (if \( a_3 = 4 \), then \( b_2 = 5 \) and \( b_1 = 2 \); and if \( a_3 = 5 \), then \( b_2 = 4 \) and \( b_1 = 2 \)).

 Case 3: \( a_2 = 4 \). At this case, \( a_3 \) must be 5. But then, the inequality \( a_2 < b_2 \) cannot be satisfied. This case is not allowed.

 Case 4: \( a_2 = 5 \). At this case, \( a_3 < a_2 \). But, we said that \( a_2 < a_3 \)! Thus, there lies a contradiction, and this case is not allowed.

 Thus, there are \( 3 + 2 = 5 \) total cases that satisfy the conditions.

 Because we don't care about the order we choose the numbers to give to \( a_1 \), \( a_2 \), and \( a_3 \) (as we could just permute them such that \( a_1 < a_2 < a_3 \)), and the same for \( b_1 \), \( b_2 \), and \( b_3 \), there are a total of \( \binom{6}{3} \cdot \binom{3}{3} = 20 \cdot 1 = 20 \) ways to distribute the elements of the set to each \( a_1 \), \( a_2 \), \( a_3 \), \( b_1 \), \( b_2 \), \( b_3 \). Thus, our probability is \( \frac{5}{20} = 1/4 \), in which our answer is \( 1+4 = \boxed{005} \).

 


	\item (\textit{AIME I 1993}) Let $S\,$ be a set with six elements. In how many different ways can one select two not necessarily distinct subsets of $S\,$ so that the union of the two subsets is $S\,$?  The order of selection does not matter; for example, the pair of subsets $\{a, c\},\{b, c, d, e, f\}$ represents the same selection as the pair $\{b, c, d, e, f\},\{a, c\}.$

 \textbf{Answer: }365



\textbf{Solution 1 (Overcounting)}

Call the two subsets $m$ and $n.$ For each of the elements in $S,$ we can assign it to either $m,n,$ or both. This gives us $3^6$ possible methods of selection. However, because the order of the subsets does not matter, each possible selection is double counted, except the case where both $m$ and $n$ contain all $6$ elements of $S.$ So our final answer is then $\frac {3^6 - 1}{2} + 1 = \boxed{365}.$

 Note: because we can order selections $A$ and $B$ in two ways, but since order does not matter in the problem, we should only count them once, ie we need to divide by two. We first need to subtract one for the count $3^6-1$ because there is one case in which $A$ and $B$ are identical and so while the rest of the “different” cases are counted twice, this one was only counted once so we should subtract one from the total count, divide by two, and add back that distinct one case. Or we could just add one, making all of the problem’s different cases counted twice, and divide it all by two, which is $\frac{3^6+1}{2}=365.$




\textbf{Solution 2 (Overcounting)}

Given one of ${6 \choose k}$ subsets with $k$ elements, the other also has $2^k$ possibilities; this is because it must contain all of the "missing" $n - k$ elements and thus has a choice over the remaining $k.$ We want $\sum_{k = 0}^6 {6 \choose k}2^k = (2 + 1)^6 = 729$ by the Binomial Theorem. But the order of the sets doesn't matter, so we get $\dfrac{729 - 1}{2} + 1 = \boxed{365}.$




\textbf{Solution 3 (Recursion)}

For all nonnegative integers $n,$ let $n$ be the number of elements in $S,$ and $f(n)$ be the number of unordered pairs $\{A,B\}$ of subsets of $S$ for which $A\cup B=S.$ We wish to find $f(6).$

 Without the loss of generality, let the elements of $S$ be $1,2,\ldots,n.$ Based on the value of $n,$ we construct the following table:
 \[\begin{array}{c|c|c} & & \\ [-2.5ex] \boldsymbol{n} & \textbf{Unordered Pairs }\boldsymbol{\{A,B\}}\textbf{ of subsets of }\boldsymbol{S}\textbf{ for which }\boldsymbol{A\cup B=S} & \boldsymbol{f(n)} \\ [0.5ex] \hline & & \\ [-2ex] 0 & \{\varnothing,\varnothing\} & 1 \\ 1 & \{\{1\},\{1\}\}, \ \{\{1\},\varnothing\} & 2 \\ 2 & \color{red}\{\{1,2\},\{1,2\}\}\color{black}, \ \color{green}\{\{1,2\},\{1\}\}\color{black}, \ \color{blue}\{\{1,2\},\{2\}\}\color{black}, \ \color{cyan}\{\{1,2\},\varnothing\}\color{black}, \ \color{magenta}\{\{1\},\{2\}\} & 5 \\ 3 & \color{red}\{\{1,2,3\},\{1,2,3\}\}\color{black}, \ \color{red}\{\{1,2,3\},\{1,2\}\}\color{black}, \ \color{green}\{\{1,2,3\},\{1,3\}\}\color{black}, \ \color{green}\{\{1,2,3\},\{1\}\}\color{black}, \ \color{green}\{\{1,2\},\{1,3\}\}\color{black}, & 14 \\ & \color{blue}\{\{1,2,3\},\{2,3\}\}\color{black}, \ \color{blue}\{\{1,2,3\},\{2\}\}\color{black}, \ \color{blue}\{\{1,2\},\{2,3\}\}\color{black}, \ \color{cyan}\{\{1,2,3\},\{3\}\}\color{black}, \ \color{cyan}\{\{1,2,3\},\varnothing\}\color{black}, \ \color{cyan}\{\{1,2\},\{3\}\}\color{black}, & \\ & \color{magenta}\{\{1,3\},\{2,3\}\}\color{black}, \ \color{magenta}\{\{1,3\},\{2\}\}\color{black}, \ \color{magenta}\{\{1\},\{2,3\}\} & \\ \vdots & \vdots & \vdots  \end{array}\]Note that for all $n\geq1,$ each unordered pair of the $(n-1)$-element set contributes three unordered pairs to the $n$-element set, as the element $n$ can be added to either or both of the subsets. The only exception is that the unordered pair of identical subsets of the $(n-1)$-element set, namely $\{\{1,2,\ldots,n-1\},\{1,2,\ldots,n-1\}\},$ only contributes two unordered pairs to the $n$-element set. The table above shows this relationship between the $2$-element set and the $3$-element set.

 Together, the recursive formula for $f(n)$ is
 \[f(n) = \begin{cases} 1 & \mathrm{if} \ n=0 \\ 3f(n-1)-1 & \mathrm{if} \ n\geq1 \end{cases}.\]Two solutions follow from here:

 



\textbf{Solution 3.1 (Recursive Formula)}

We evaluate $f(6)$ recursively:
 \begin{alignat*}{6} f(0)&=1, \\ f(1)&=3f(0)-1&&=2, \\ f(2)&=3f(1)-1&&=5, \\ f(3)&=3f(2)-1&&=14, \\ f(4)&=3f(3)-1&&=41, \\ f(5)&=3f(4)-1&&=122, \\ f(6)&=3f(5)-1&&=\boxed{365}. \end{alignat*}~MRENTHUSIASM

 



\textbf{Solution 3.2 (Explicit Formula)}

For all $n\geq1,$ we have
 \begin{align*} f(n) &= 3f(n-1)-1 \\ &= 3\left(3f(n-2)-1\right)-1 \\ &= 3^2f(n-2)-3-1 \\ &= 3^2\left(3f(n-3)-1\right)-3-1 \\ &= 3^3f(n-3)-3^2-3-1 \\ & \ \vdots \\ &= 3^nf(0)-3^{n-1}-3^{n-2}-3^{n-3}-\cdots-1 \\ &= 3^n-\left(3^{n-1}+3^{n-2}+3^{n-3}+\cdots+1\right) \\ &= 3^n-\frac{3^n-1}{2} \\ &= \frac{3^n+1}{2} \\ &= \frac{3^n-1}{2}+1, \end{align*}
which resembles the result in Solutions 1 and 2. The answer is $f(6)=\boxed{365}.$

 As $f(0)=1,$ note that 
 \[f(n)=\frac{3^n-1}{2}+1\] holds for $n=0$ too. Therefore, this formula is true for all nonnegative integers $n.$




\textbf{Solution 4 (Casework)}

We can perform casework based on the number of overlapping elements. If no elements overlap, there is $\binom60=1$ way to choose the overlapping elements, and $2^{6-0}$ ways to distribute the remaining elements--each element can go in one subset or the other. We must also divide by $2$ because the order of the subsets does not matter. Proceeding similarly for the other cases, our sum is 
 \[\dbinom{6}0 \cdot 2^5+\dbinom{6}1 \cdot 2^4+\cdots +\dbinom{6}4 \cdot 2+\dbinom{6}5 \cdot 1+\dbinom{6}6 \cdot 1.\](Note that we have to be especially careful with the last case, as it does not follow the pattern of the other cases.) Adding these up gives a total of $\boxed{365}.$




\textbf{Solution 5 (Bash)}

If we wanted to, we could use casework, being very careful not to double count cases. Let's first figure out what cases we need to look at. The notation I will be using is $x\rightarrow y,$ which implies that we pick $x$ numbers for the first set which then the second set can have $y$ numbers.

 Clearly:
 \begin{align*} 0&\rightarrow6 \\ 1&\rightarrow5\mid6 \\ 2&\rightarrow4\mid5\mid6 \\ 3&\rightarrow3\mid4\mid5\mid6 \\ 4&\rightarrow2\mid3\mid4\mid5\mid6 \\ 5&\rightarrow1\mid2\mid3\mid4\mid5\mid6 \\ 6&\rightarrow0\mid1\mid2\mid3\mid4\mid5\mid6 \end{align*}
However notice that many of the cases are double counted as direction does not matter, e.g. $2\rightarrow4$ is the same as $4\rightarrow2.$ Get rid of those cases so we are just left with:
 \begin{align*} 0&\rightarrow6 \\ 1&\rightarrow5\mid6 \\ 2&\rightarrow4\mid5\mid6 \\ 3&\rightarrow3\mid4\mid5\mid6 \\ 4&\rightarrow4\mid5\mid6 \\ 5&\rightarrow5\mid6 \\ 6&\rightarrow6 \end{align*}
Now we start computing, $0\rightarrow6$ is just $1$ case, $1\rightarrow5\mid6$ is just $\binom{6}{1}\cdot2$ cases, $2\rightarrow4\mid5\mid6$ is just $\binom{6}{2}\cdot2^2$ cases, and $3\rightarrow3\mid4\mid5\mid6$ is just $\binom{6}{3}\cdot2^3$ cases (If you have trouble understanding this, write down the six letters and then try to understand what $x\rightarrow y$ really means.). Now what you can do is continue on this same pattern like Solution 2 does, and then use simple symmetry to figure out the double counted cases. However, the purpose of this solution is to bash out the double counted cases, so we will do exactly that.

 One quick thing though. We have a double counted case with the $3\rightarrow3,$ as choosing ABC and DEF is the same thing as choosing DEF and then ABC. There are $\frac{\binom{6}{3}}{2} = 10$ cases of this.

 For computing $4\rightarrow4\mid5\mid6,$ we use the same process as before. We have $\binom{6}{4}\cdot(3+4+1)$ (Note, the $3$ comes from $\frac{\binom{4}{2}}{2}$), and for $5\rightarrow5\mid6$ we have $\binom{6}{5}\cdot\left(\frac{5}{2}+1\right),$ and then for $6\rightarrow6$ we just have $\binom{6}{6}$ (there is no double counted case since ABCDEF, ABCDEF is only counted once).

 Summing case by case, we have $1+12+60+150+120+21+1 = \boxed{365}.$



	\item (\textit{AIME I 1993}) Two thousand points are given on a circle. Label one of the points 1. From this point, count 2 points in the clockwise direction and label this point 2. From the point labeled 2, count 3 points in the clockwise direction and label this point 3. (See figure.) Continue this process until the labels $1,2,3\dots,1993\,$ are all used. Some of the points on the circle will have more than one label and some points will not have a label. What is the smallest integer that labels the same point as 1993? 


 
                    \begin{figure}[H]
                        \centering
                        \includegraphics[width=0.5\textwidth]{Problem/AIME_I_1993_9_Problem_Figure_1.png}
                    \end{figure}
                    

 \textbf{Answer: }118



\textbf{Solution}

The label $1993$ will occur on the $\frac12(1993)(1994) \pmod{2000}$th point around the circle. (Starting from 1) A number $n$ will only occupy the same point on the circle if $\frac12(n)(n + 1)\equiv \frac12(1993)(1994) \pmod{2000}$. 

 Simplifying this expression, we see that $(1993)(1994) - (n)(n + 1) = (1993 - n)(1994 + n)\equiv 0\pmod{2000}$. Therefore, one of $1993 - n$ or $1994 + n$ is odd, and each of them must be a multiple of $125$ or $16$. 

 For $1993 - n$ to be a multiple of $125$ and $1994 + n$ to be a multiple of $16$, $n \equiv 118 \pmod {125}$ and $n\equiv 6 \pmod {16}$. The smallest $n$ for this case is $118$. 

 In order for $1993 - n$ to be a multiple of $16$ and $1994 + n$ to be a multiple of $125$, $n\equiv 9\pmod{16}$ and $n\equiv 6\pmod{125}$. The smallest $n$ for this case is larger than $118$, so $\boxed{118}$ is our answer.

 \textbf{Note:} One can just substitute $1993\equiv-7\pmod{2000}$ and $1994\equiv-6\pmod{2000}$ to simplify calculations.

 



\textbf{Solution 2}

Two labels $a$ and $b$ occur on the same point if $\ a(a+1)/2\equiv \ b(b+1)/2\pmod{2000}$. If we assume the final answer be $n$, then we have $\frac12(n)(n + 1)\equiv \frac12(1993)(1994) \pmod{2000}$.

 Multiply $2$ on both side we have $(1993)(1994) - (n)(n + 1) = (1993 - n)(1994 + n)\equiv 0\pmod{4000}$. As they have different parities, the even one must be divisible by $32$. As $(1993 - n)+(1994 + n)\equiv 2\pmod{5}$, one of them is divisible by $5$, which indicates it's divisible by $125$.

 Which leads to four different cases: $1993-n\equiv 0\pmod{4000}$ ; $1994+n\equiv 0\pmod{4000}$ ; $1993-n\equiv 0\pmod{32}$ and $1994+n\equiv 0\pmod{125}$ ; $1993-n\equiv 0\pmod{125}$ and $1994+n\equiv 0\pmod{32}$. Which leads to $n\equiv 1993,2006,3881$ and $118\pmod{4000}$ respectively, and only $n=118$ satisfied.Therefore answer is $\boxed{118}$.(by ZJY)

 


	\item (\textit{AIME I 1993}) Euler's formula states that for a convex polyhedron with $V\,$ vertices, $E\,$ edges, and $F\,$ faces, $V-E+F=2\,$. A particular convex polyhedron has 32 faces, each of which is either a triangle or a pentagon. At each of its $V\,$ vertices, $T\,$ triangular faces and $P^{}_{}$ pentagonal faces meet. What is the value of $100P+10T+V\,$? 


 
 \textbf{Answer: }250



\textbf{Solution 1}

The convex polyhedron of the problem can be easily visualized; it corresponds to a \href{https://artofproblemsolving.com/wiki/index.php?title=Dodecahedron}{dodecahedron} (a regular solid with $12$ \href{https://artofproblemsolving.com/wiki/index.php?title=Equilateral}{equilateral} pentagons) in which the $20$ vertices have all been truncated to form $20$ equilateral triangles with common vertices. The resulting solid has then $p=12$ smaller equilateral pentagons and $t=20$ equilateral triangles yielding a total of $t+p=F=32$ faces. In each vertex, $T=2$ triangles and $P=2$ pentagons are concurrent. Now, the number of edges $E$ can be obtained if we count the number of sides that each triangle and pentagon contributes: $E=\frac{3t+5p}{2}$, (the factor $2$ in the \href{https://artofproblemsolving.com/wiki/index.php?title=Denominator}{denominator} is because we are counting twice each edge, since two adjacent faces share one edge). Thus, $E=60$. Finally, using Euler's formula we have $V=E-30=30$.

 In summary, the solution to the problem is $100P+10T+V=\boxed{250}$.

 



\textbf{Solution 2}

As seen above, $E=V+30$. Every vertex $V$, there is a triangle for every $T$ and a pentagon for every $P$ by the given. However, there are three times every triangle will be counted and five times every pentagon will be counted because of their numbers of vertices. From this observation, $\frac{VT}3+\frac{VP}5=32\implies V(5T+3P)=480$. Also, at every vertex $V$, there are $T+P$ edges coming out from that vertex (one way to see this is to imagine the leftmost segment of each triangle and pentagon that is connected to the given vertex, and note that it includes every one of the edges exactly once), so $\frac{V(T+P)}2=E\implies V(T+P)=2E\implies V(5T+5P)=10E$, and subtracting the other equation involving the vertices from this gives$2VP=10E-480\implies VP=5E-240=5(V+30)-240=5V-90$$\implies V(5-P)=90$. Since $V|480$ from the first vertex-related observation and $P>0\implies5-P<5$, and it quickly follows that $V=30\implies E=60\implies P=2\implies T=2\implies100P+10T+V=\boxed{250}$.

 


 



\textbf{Solution 3}

Notice that at each human, we must have the sum of the angles be less than $360$ degrees or we will not be able to fold the polyhedron.  Therefore, we have $60T + 108P < 360.$ Now, let there be $t$ triangles and $p$ pentagons total such that $t+p = 32.$  From the given, we know that $E = V + 30.$  Lastly, we see that $E = \frac{3t+5p}{2}$ and $V = \frac{3t}{T}=\frac{5p}{P}.$ 

 Now, we do casework on what $P$ is.  

 Case 1: $P = 2$Notice that we must have $t$ and $p$ integral.  Trying $T = 1, 2$ yields a solution with $t=2.$  Trying other cases of $P$ and $T$ yields no solutions.  Therefore, $T=2, P=2$ and after solving for $t, p$ we get $V=30.$  Finally, we have $100P+10T+V = \boxed{250}$.

 



\textbf{Solution 4}

We know that $V-E = -30 \implies V = E-30$ based off the problem condition. Furthermore, if we draw out a few pentagons as well as triangles on each of side of the pentagons, it's clear that each vertex has 4 edges connected to it, with two triangles and two pentagons for each vertex. However, each edge is used for two vertices and thus counted twice. Therefore, we have that $E = \frac{4V}{2} = 2V$. Plugging this in, we have that $V=30$ and so our answer is $200+20+30 = \boxed{250}$.

 \textbf{Note} The solid described in this problem is an icosidodecahedron.

 


	\item (\textit{AIME I 1993}) Alfred and Bonnie play a game in which they take turns tossing a fair coin. The winner of a game is the first person to obtain a head. Alfred and Bonnie play this game several times with the stipulation that the loser of a game goes first in the next game. Suppose that Alfred goes first in the first game, and that the probability that he wins the sixth game is $m/n\,$, where $m\,$ and $n\,$ are relatively prime positive integers. What are the last three digits of $m+n\,$? 


 
 \textbf{Answer: }093



\textbf{Solution}

The probability that the $n$th flip in each game occurs and is a head is $\frac{1}{2^n}$. The first person wins if the coin lands heads on an odd numbered flip. So, the probability of the first person winning the game is $\frac{1}{2}+\frac{1}{8}+\frac{1}{32}+\cdots = \frac{\frac{1}{2}}{1-\frac{1}{4}}=\frac{2}{3}$, and the probability of the second person winning is $\frac{1}{3}$. 

 Let $a_n$ be the probability that Alfred wins the $n$th game, and let $b_n$ be the probability that Bonnie wins the $n$th game. 

 If Alfred wins the $n$th game, then the probability that Alfred wins the $n+1$th game is $\frac{1}{3}$. If Bonnie wins the $n$th game, then the probability that Alfred wins the $n+1$th game is $\frac{2}{3}$. 

 Thus, $a_{n+1}=\frac{1}{3}a_n+\frac{2}{3}b_n$. 

 Similarly, $b_{n+1}=\frac{2}{3}a_n+\frac{1}{3}b_n$. 

 Since Alfred goes first in the $1$st game, $(a_1,b_1)=\left(\frac{2}{3}, \frac{1}{3}\right)$. 

 Using these \href{https://artofproblemsolving.com/wiki/index.php?title=Recursion}{recursive} equations: 

 $(a_2,b_2)=\left(\frac{4}{9}, \frac{5}{9}\right)$
$(a_3,b_3)=\left(\frac{14}{27}, \frac{13}{27}\right)$
$(a_4,b_4)=\left(\frac{40}{81}, \frac{41}{81}\right)$
$(a_5,b_5)=\left(\frac{122}{243}, \frac{121}{243}\right)$
$(a_6,b_6)=\left(\frac{364}{729}, \frac{365}{729}\right)$ 

 Since $a_6=\frac{364}{729}$, $m+n = 1093 \equiv \boxed{093} \pmod{1000}$.

 \textbf{Simpler Solution 1} Same as solution 1 except that $b_n=1-a_n$. So you don't need that extra calculation for $b_n$.

 



\textbf{Solution 2 (Kind of Bashy but very nice)}

In order to begin this problem, we need to calculate the probability that Alfred will win on the first round.

 Because he goes first, Alfred has a $\frac{1}{2}$ chance of winning (getting heads) on his first flip.

 Then, Bonnie, who goes second, has a $\frac{1}{2} \times \frac{1}{2} = \frac{1}{4}$, chance of winning on her first coin toss. 

 Therefore Alfred’s chance of winning on his second flip is $\frac{1}{4} \times \frac{1}{2} = \frac{1}{8}$

 


 From this, we can see that Alfred’s (who goes first) chance of winning the first round is: $\frac{1}{2} + \frac{1}{8} + \frac{1}{32} + \cdots = \frac{2}{3}$

 Bonnie’s (who goes second) chance of winning the first round is then $1 - \frac{2}{3} = \frac{1}{3}$.

 This means that the person who goes first has a $\frac{2}{3}$ chance of winning the round, while the person who goes second has a $\frac{1}{3}$ chance of winning.

 


 Now, through casework, we can calculate Alfred’s chance of winning the second round.

 Case 1: Alfred wins twice; $\frac{2}{3} \times \frac{1}{3}$ (Bonnie goes first this round) $=\frac{2}{9}$.

 Case 2: Alfred loses the first round, but wins the second; $\frac{1}{3} \times \frac{2}{3} = \frac{2}{9}$.

 Adding up the cases, we get $\frac{2}{9} + \frac{2}{9} = \frac{4}{9}$.

 Alfred, therefore, has a $\frac{4}{9}$ of winnning the second round, and Bonnie has a $1-\frac{4}{9} = \frac{5}{9}$ chance of winning this round.

 


 From here, it is not difficult to see that the probabilities alternate in a pattern.

 Make $A$ the probability that Alfred wins a round.

 The chances of Alfred and Bonnie, respectively, winning the first round are $A$ and $A - \frac{1}{3}$, which can be written as $2A - \frac{1}{3} = 1$

 The second round’s for chances are $A$ and $A + \frac{1}{9}$, which can also be written as $2A + \frac{1}{9}$

 


 From this, we can conclude that for the $n$th even round, the probability that Alfred ($A$) wins can be calculated through the equation $2A + \frac{1}{3^n} = 1$.

 Solving the equation for $n = 6$, we get $A = \frac{364}{729}$.

 $364 + 729 = 1093$.

 So our answer is $\boxed{093}$.

 



\textbf{Solution 3}

Rather than categorizing games as wins or losses, we can categorize them as starters (S), where Alfred starts, and non-starters (NS), where Bonnie starts. Game 1 is a starter, and since Alfred must win Game 6, Game 7 is a non-starter.

 As shown in Solution 1, if a player starts a certain game, the probability $P(NS)$ that they will not start the next game (i.e. win the current game) is $\frac{2}{3}$, and the probability $P(S)$ that they will start the next game (i.e. lose the current game) is $\frac{1}{3}$. Similarly, if a player does not start a certain game, $P(NS) = \frac{1}{3}$ and $P(S) = \frac{2}{3}$. We conclude that the probability of switching from S to NS or vice versa is always $\frac{2}{3}$, and the probability of staying the same is always $\frac{1}{3}$.

 Listing out all the games from Game 1 to Game 7, we get: S, ?, ?, ?, ?, ?, NS. Between the 7 games, there are 6 opportunities to either switch or stay the same. We need to eventually switch from S to NS, so there must be an odd number of switches. Furthermore, this number is less than 6, so it must be 1, 3, or 5.

 1 switch: There are ${6 \choose 1} = 6$ ways to place the switch, so $P = 6\left(\frac{1}{3}\right)^5\left(\frac{2}{3}\right) = \frac{12}{729}$.

 3 switches: There are ${6 \choose 3} = 20$ ways to place the switches, so $P = 20\left(\frac{1}{3}\right)^3\left(\frac{2}{3}\right)^3 = \frac{160}{729}$.

 5 switches: There are ${6 \choose 5} = 6$ ways to place the switch, so $P = 6\left(\frac{1}{3}\right)\left(\frac{2}{3}\right)^5 = \frac{192}{729}$.

 Add up all these probabilities to get $\frac{12+160+192}{729} = \frac{364}{729}$.   $364+729=1093$, so the answer is $\boxed{093}$.

 



\textbf{Solution 4 (Probability States)}

In states problems like this, it is always good to find a relation between the first game and relate it to all other games.

 To do such a thing, we first develop a Markov Chain, shown below.

 
                    \begin{figure}[H]
                        \centering
                        \includegraphics[width=0.5\textwidth]{Solution/AIME_I_1993_11_Solution_Figure_1.png}
                    \end{figure}
                    Now, say \( A \) is the probability Alfred wins and \( B \) is the probability Bonnie wins, given that they each go first. Then, one can notice that \( A \) will win either immediately with a chance of \( 1/2 \), or will lose his turn with a chance \( 1/2 \). But then, one can find a bijection that the new case is just the probability Bonnie wins, given that she goes first. However, we want Alfred to win, so the requested probability \( A \) is
 \[A = \frac{1}{2} + \frac{1}{2}(1-B).\]Similarly for \( B \),
 \[B = \frac{1}{2} + \frac{1}{2}(1-A).\]We can solve for any one of \( A \) or \( B \), as they both will retain the same probability due to the fact they both start, and we can clearly see that solving for \( B \);
 \[B = \frac{1}{2} + \frac{1}{2}(1-A) = \frac{1}{2} + \frac{1}{2} - \frac{1}{2}A = 1 - \frac{1}{2}A\]Substituting the expression for \( A \):
 \[B = 1 - \frac{1}{2}\left(\frac{1}{2} + \frac{1}{2}(1-B)\right) = 1 - \frac{1}{2}\left(1 - \frac{1}{2}B\right) = 1 - \frac{1}{2} + \frac{1}{4}B = \frac{1}{2} + \frac{1}{4}B\]in which \( B = \frac{1}{2} + \frac{1}{4}B \), or \( \frac{3}{4}B = \frac{1}{2} \), or \( B = \frac{2}{3} \). Thus, the probability that Alfred wins the first game given he goes first is \( 2/3 \).

 Now, let \( P_n \) be the probability that on the \( n \)th game, Alfred wins. One can notice that if Alfred starts, that implies he lost the last game, and thus has a \( 2/3 (1-P_{n-1}) \) chance of winning. However, one can also notice that if Alfred goes second, then he has a \( 1/3(P_{n-1}) \) chance of winning. Thus the recurrence relation
 \[P_n = \frac{2}{3} (1-P_{n-1}) + \frac{1}{3}(P_{n-1})\]is formed.

 The recurrence simplifies down to
 \[P_n = \frac{2}{3} - \frac{1}{3}(P_{n-1}) \quad : \quad P_1 = \frac{2}{3}.\]We can bash our way to \( P_6 \) quite easily (which for the sake of spacing, I will not do). One can also calculate this explicitly, to obtain \( P_6 = 364/729 \), which is in the form \( m/n \), and \( m+n = 364+729 = 1093 \). The last three digits are \( \boxed{093} \).

 


	\item (\textit{AIME I 1993}) The vertices of $\triangle ABC$ are $A = (0,0)\,$, $B = (0,420)\,$, and $C = (560,0)\,$.  The six faces of a die are labeled with two $A\,$'s, two $B\,$'s, and two $C\,$'s.  Point $P_1 = (k,m)\,$ is chosen in the interior of $\triangle ABC$, and points $P_2\,$, $P_3\,$, $P_4, \dots$ are generated by rolling the die repeatedly and applying the rule: If the die shows label $L\,$, where $L \in \{A, B, C\}$, and $P_n\,$ is the most recently obtained point, then $P_{n + 1}^{}$ is the midpoint of $\overline{P_n L}$.  Given that $P_7 = (14,92)\,$, what is $k + m\,$?


 
 \textbf{Answer: }344



\textbf{Solution 1}

If we have points $(p,q)$ and $(r,s)$ and we want to find $(u,v)$ so $(r,s)$ is the midpoint of $(u,v)$ and $(p,q)$, then $u=2r-p$ and $v=2s-q$. So we start with the point they gave us and work backwards. We make sure all the coordinates stay within the triangle. We have 
 \[P_{n-1}=(x_{n-1},y_{n-1}) = (2x_n\bmod{560},\ 2y_n\bmod{420})\]Then $P_7=(14,92)$, so $x_7=14$ and $y_7=92$, and we get 
 \[\begin{array}{c||ccccccc} n & 7 & 6 & 5 & 4 & 3 & 2 & 1 \\ \hline\hline x_n & 14 & 28 & 56 & 112 & 224 & 448 & 336 \\ \hline y_n & 92 & 184 & 368 & 316 & 212 & 4 & 8  \end{array}\]

 So the answer is $\boxed{344}$.

 



\textbf{Solution 2}

Let $L_1$ be the $n^{th}$ roll that directly influences $P_{n + 1}$. 

 Note that $P_7 = \cfrac{\cfrac{\cfrac{P_1 + L_1}2 + L_2}2 + \cdots}{2\ldots} = \frac {(k,m)}{64} + \frac {L_1}{64} + \frac {L_2}{32} + \frac {L_3}{16} + \frac {L_4}8 + \frac {L_5}4 + \frac {L_6}2 = (14,92)$. 

 
Then quickly checking each addend from the right to the left, we have the following information (remembering that if a point must be $(0,0)$, we can just ignore it!): 

 for $\frac {L_6}2,\frac {L_5}4$, since all addends are nonnegative, a non-$(0,0)$ value will result in a $x$ or $y$ value greater than $14$ or $92$, respectively, and we can ignore them,

 for $\frac {L_4}8,\frac {L_3}{16},\frac {L_2}{32}$ in a similar way, $(0,0)$ and $(0,420)$ are the only possibilities, 

 and for $\frac {L_1}{64}$, all three work. 

 


 Also, to be in the triangle, $0\le k\le560$ and $0\le m\le420$.

 
Since $L_1$ is the only point that can possibly influence the $x$ coordinate other than $P_1$, we look at that first.

 If $L_1 = (0,0)$, then $k = 2^6\cdot14 = 64\cdot14 > 40\cdot14 = 560$, 

 so it can only be that $L_1 = (560,0)$, and $k + 560 = 2^6\cdot14$
$\implies k = 64\cdot14 - 40\cdot14 = 24\cdot14 = 6\cdot56 = 336$.

 
Now, considering the $y$ coordinate, note that if any of $L_2,L_3,L_4$ are $(0,0)$ ($L_2$ would influence the least, so we test that), 

 then $\frac {L_2}{32} + \frac {L_3}{16} + \frac {L_4}8 < \frac {420}{16} + \frac {420}8 = 79\pm\epsilon < 80$, 

 which would mean that $P_1 > 2^6\cdot(92 - 80) = 64\cdot12 > 42\cdot10 = 420\ge m$, so $L_2,L_3,L_4 = (0,420)$, 

 and now $\frac {P_1}{64} + \frac {420}{2^5} + \frac {420}{2^4} + \frac {420}{2^3} = 92$
$\implies P_1$
$= 64\cdot92 - 420(2 + 4 + 8)$
$= 64\cdot92 - 420\cdot14= 64(100 - 8) - 14^2\cdot30$
$= 6400 - 512 - (200 - 4)\cdot30$
$= 6400 - 512 - 6000 + 120$
$= - 112 + 120$
$= 8$, 

 and finally, $k + m = 336 + 8 = \boxed{344}$.

 


	\item (\textit{AIME I 1993}) Jenny and Kenny are walking in the same direction, Kenny at 3 feet per second and Jenny at 1 foot per second, on parallel paths that are 200 feet apart. A tall circular building 100 feet in diameter is centered midway between the paths. At the instant when the building first blocks the line of sight between Jenny and Kenny, they are 200 feet apart. Let $t\,$ be the amount of time, in seconds, before Jenny and Kenny can see each other again. If $t\,$ is written as a fraction in lowest terms, what is the sum of the numerator and denominator?


 
 \textbf{Answer: }163



\textbf{Solution 1}

Consider the unit cicle of radius 50. Assume that they start at points $(-50,100)$ and $(-50,-100).$ Then at time $t$, they end up at points $(-50+t,100)$ and $(-50+3t,-100).$The equation of the line connecting these points and the equation of the circle are 
 \begin{align}y&=-\frac{100}{t}x+200-\frac{5000}{t}\\50^2&=x^2+y^2\end{align}. When they see each other again, the line connecting the two points will be tangent to the circle at the point $(x,y).$ Since the radius is perpendicular to the tangent we get 
 \[-\frac{x}{y}=-\frac{100}{t}\] or $xt=100y.$ Now substitute 
 \[y= \frac{xt}{100}\] into $(2)$ and get 
 \[x=\frac{5000}{\sqrt{100^2+t^2}}.\] Now substitute this and 
 \[y=\frac{xt}{100}\] into $(1)$ and solve for $t$ to get 
 \[t=\frac{160}{3}.\] Finally, the sum of the numerator and denominator is $160+3=\boxed{163}.$




\textbf{Solution 2}

Let $A$ and $B$ be Kenny's initial and final points respectively and define $C$ and $D$ similarly for Jenny. Let $O$ be the center of the building. Also, let $X$ be the intersection of $AC$ and $BD$. Finaly, let $P$ and $Q$ be the points of tangency of circle $O$ to $AC$ and $BD$ respectively.  

 
                    \begin{figure}[H]
                        \centering
                        \includegraphics[width=0.5\textwidth]{Solution/AIME_I_1993_13_Solution_Figure_1.png}
                    \end{figure}
                    

 From the problem statement, $AB=3t$, and $CD=t$. Since $\Delta ABX \sim \Delta CDX$, $CX=AC\cdot\left(\frac{CD}{AB-CD}\right)=200\cdot\left(\frac{t}{3t-t}\right)=100$. 

 Since $PC=100$, $PX=200$. So, $\tan(\angle OXP)=\frac{OP}{PX}=\frac{50}{200}=\frac{1}{4}$. 

 Since circle $O$ is tangent to $BX$ and $AX$, $OX$ is the \href{https://artofproblemsolving.com/wiki/index.php?title=Angle\_bisector}{angle bisector} of $\angle BXA$. 

 Thus, $\tan(\angle BXA)=\tan(2\angle OXP)=\frac{2\tan(\angle OXP)}{1- \tan^2(\angle OXP)} = \frac{2\cdot \left(\frac{1}{4}\right)}{1-\left(\frac{1}{4}\right)^2}=\frac{8}{15}$. 

 Therefore, $t = CD = CX\cdot\tan(\angle BXA) = 100 \cdot \frac{8}{15} = \frac{160}{3}$, and the answer is $\boxed{163}$.

 



\textbf{Solution 3}


                    \begin{figure}[H]
                        \centering
                        \includegraphics[width=0.5\textwidth]{Solution/AIME_I_1993_13_Solution_Figure_2.png}
                    \end{figure}
                    

 Let $t$ be the time they walk.  Then $CD=t$ and $AB=3t$.

 
Draw a line from point $O$ to $Q$ such that $OQ$ is perpendicular to $BD$.  Further, draw a line passing through points $O$ and $P$, so $OP$ is parallel to $AB$ and $CD$ and is midway between those two lines.  Then $PR=\dfrac{AB+CD}{2}=\dfrac{3t+t}{2}=2t$.  Draw another line passing through point $D$ and parallel to $AC$, and call the point of intersection of this line with $AB$ as $S$.  Then $SB=AB-AS=3t-t=2t$.

 
We see that $m\angle SBD=m\angle ORQ$ since they are corresponding angles, and thus by angle-angle similarity, $\triangle QOR\sim\triangle SDB$.

 
Then
 \begin{align*} \dfrac{OQ}{DS}=\dfrac{RO}{BD}&\implies\dfrac{50}{200}=\dfrac{RO}{\sqrt{200^2+4t^2}}\\ &\implies RO=\dfrac{1}{4}\left(\sqrt{200^2+4t^2}\right)\\ &\implies RO=\dfrac{1}{2}\left(\sqrt{100^2+t^2}\right) \end{align*}


 
And we obtain
 \begin{align*} PR-OP&=RO\\ 2t-50&=\dfrac{1}{2}\left(\sqrt{100^2+t^2}\right)\\ 4t-100&=\sqrt{100^2+t^2}\\ (4t-100)^2&=\left(\sqrt{100^2+t^2}\right)^2\\ 16t^2-800t+100^2&=t^2+100^2\\ 15t^2&=800t\\ t&=\dfrac{800}{15} \end{align*}


 so we have $t=\frac{160}{3}$, and our answer is thus $160+3=\boxed{163}$.

 



\textbf{Solution 4}

We can use areas to find the answer. Since Jenny and Kenny are 200 feet apart, we know that they are side by side, and that the line connecting the two of them is tangent to the circular building. 

 We know that the radius of the circle is 50, and that $\overline{AJ} = x$, $\overline{BK} = 3x$.

 Illustration:

 
                    \begin{figure}[H]
                        \centering
                        \includegraphics[width=0.5\textwidth]{Solution/AIME_I_1993_13_Solution_Figure_3.png}
                    \end{figure}
                    

 By areas, $[OJK] + [AJOF] + [OFBK] = [ABKJ]$. Having right trapezoids, $[AFOJ] = \frac{x+50}{2} \cdot 100$. The other areas of right trapezoids can be calculated in the same way. We just need to find $[OJK]$ in terms of $x$.

 If we bring $\overline{AB}$ up to where the point J is, we have by the Pythagorean Theorem, $\overline{JK} = 2\sqrt{x^2+10000} \Rightarrow [OJK] = \frac{1}{2} \overline{JK} \cdot \overline{OD}$.

 Now we have everything to solve for $x$.

 
 \[[OJK] + [AJOF] + [OFBK] = [ABKJ]\]

 \[50 \sqrt{x^2 + 10000} + \frac{x+50}{2} \cdot 100 + \frac{3x+50}{2} \cdot 100 = \frac{x+3x}{2} \cdot 200\]

 After isolating the radical, dividing by 50, and squaring, we obtain: $15x^2 - 800x = 0 \Rightarrow x = \frac{160}{3}$.

 Since Jenny walks $\frac{160}{3}$ feet at 1 foot/s, our answer is $160+3 = \fbox{163}$.

 



\textbf{Solution 5}


                    \begin{figure}[H]
                        \centering
                        \includegraphics[width=0.5\textwidth]{Solution/AIME_I_1993_13_Solution_Figure_4.png}
                    \end{figure}
                    Consider our diagram here, where Jenny goes from $A$ to $D$ and Kenny goes from $B$ to $C$. Really what we are asking is the length of $AD$, knowing that $AB$ and $CD$ are tangents and $AB$ is perpendicular to the parallel lines. Draw lines $EF$ and $GH$ as shown such that they are parallel to $BC$ and $AD$ and are tangent to the circle. Additionally, let $MN$ be the median of trapezoid $ABCD$. Notice how $AH=MH=50$ and $BE=ME=50$ as well, so $EF$ and $GH$ are medians of trapezoids $MBCN$ and $AMND$, respectively. If we set $AD=x$ and $BC=3x$, then $MN=2x$ so $HG=\frac{3}{2}x$ and $EF=\frac{5}{2}x$.

 Now we will find $FG$ in two different ways. For the first way, notice that $EFGH$ has its inscribed circle, so $EF+GH=EH+FG$. This means that $FG=\frac{3}{2}x+\frac{5}{2}x-100=4x-100$. However, if we let $P$ be the altitude from $G$ to $EF$, then notice how $PF=x$. Since $PG=100$, then by the Pythagorean Theorem, we also have that $FG=\sqrt{x^2+10000}$. After that, then we have an equation in $x$, which is 
 \[\sqrt{x^2+10000}=4x-100.\] Squaring both sides, it expands to 
 \[x^2+10000=16x^2-800x+10000,\] and since $x$ is nonzero, then once we cancel $10000$ from both sides, we can divide $x$ and get the new equation 
 \[x=16x-800\implies x=\frac{160}{3}.\] Because Jenny travels this at $1$ foot per second, then it will take $\frac{160}{3}$ seconds, so our answer is $\boxed{163}$.

 



\textbf{Solution 6}

Let the point where Jenny started be $A = (0,0)$. Then let $D = (t,0)$ the point Jenny is when she can see Kenny. Similarly, let $B = (0,200)$ be where Kenny started, and let $C = (3t,200)$ the point Kenny is when she can see Jenny. Then the equation of the line passing through $C, D$ is $100x-ty-100t=0$. Now, note that the center of the building is at $O = (50,100)$, and the distance from $O$ to $CD$ is $50$, so, using the distance from a point to a line formula, we get 
 \[\frac{|100\cdot50-100a-100a|}{\sqrt{10000+a^2}} = 50,\]which simplifies to $15a^2-800a = 0$ from clearing denominators, squaring to remove absolute value and radical, then cancelling and rearranging. But this is equivalent to $a = \frac{160}{3}$, so our answer is $\boxed{163}$.

 


	\item (\textit{AIME I 1993}) A rectangle that is inscribed in a larger rectangle (with one vertex on each side) is called \textit{unstuck} if it is possible to rotate (however slightly) the smaller rectangle about its center within the confines of the larger. Of all the rectangles that can be inscribed unstuck in a 6 by 8 rectangle, the smallest perimeter has the form $\sqrt{N}\,$, for a positive integer $N\,$. Find $N\,$.


 
 \textbf{Answer: }448



\textbf{Solution 1}

Put the rectangle on the coordinate plane so its vertices are at $(\pm4,\pm3)$, for all four combinations of positive and negative. Then by symmetry, the other rectangle is also centered at the origin, $O$.

 Note that such a rectangle is unstuck if its four vertices are in or on the edge of all four quadrants, and it is not the same rectangle as the big one. Let the four vertices of this rectangle be $A(4,y)$, $B(-x,3)$, $C(-4,-y)$ and $D(x,-3)$ for nonnegative $x,y$. Then this is a rectangle, so $OA=OB$, or $16+y^2=9+x^2$, so $x^2=y^2+7$.

 
                    \begin{figure}[H]
                        \centering
                        \includegraphics[width=0.5\textwidth]{Solution/AIME_I_1993_14_Solution_Figure_1.png}
                    \end{figure}
                    Reflect $D$ across the side of the rectangle containing $C$ to $D'(-8-x,-3)$. Then $BD'=\sqrt{(-8-x-(-x))^2+(3-(-3))^2}=10$ is constant, and the perimeter of the rectangle is equal to $2(BC+CD')$. The midpoint of $\overline{BD'}$ is $(-4-x,0)$, and since $-4>-4-x$ and $-y\le0$, $C$ always lies below $\overline{BD'}$.

 If $y$ is positive, it can be decreased to $y'<y$. This causes $x$ to decrease as well, to $x'$, where $x'^2=y'^2+7$ and $x'$ is still positive. If $B$ and $D'$ are held in place as everything else moves, then $C$ moves $(y-y')$ units up and $(x-x')$ units left to $C'$, which must lie within $\triangle BCD'$. Then we must have $BC'+C'D'<BC+CD'$, and the perimeter of the rectangle is decreased. Therefore, the minimum perimeter must occur with $y=0$, so $x=\sqrt7$.

 By the distance formula, this minimum perimeter is 
 \[2\left(\sqrt{(4-\sqrt7)^2+3^2}+\sqrt{(4+\sqrt7)^2+3^2}\right)=4\left(\sqrt{8-2\sqrt7}+\sqrt{8+2\sqrt7}\right)\] 
 \[=4(\sqrt7-1+\sqrt7+1)=8\sqrt7=\sqrt{448}.\] Therefore $N$ would equal $\boxed{448}.$




\textbf{Solution 2}

Note that the diagonal of the rectangle with minimum perimeter must have the diagonal along the middle segment of length 8 of the rectangle (any other inscribed rectangle can be rotated a bit, then made smaller; this one can't because then the rectangle cannot be inscribed since its longest diagonal is less than 8 in length). Then since a rectangle must have right angles, we draw a circle of radius 4 around the center of the rectangle. Picking the two midpoints on the sides of length 6 and opposite intersection points on the segments of length 8, we form a rectangle. Let $a$ and $b$ be the sides of the rectangle. Then $ab = 3(8) = 24$ since both are twice the area of the same right triangle, and $a^2+b^2 = 64$. So $(a+b)^2 = 64+2(24) = 112$, so $2(a+b) = \sqrt{\boxed{448}}$.

 


	\item (\textit{AIME I 1993}) Let $\overline{CH}$ be an altitude of $\triangle ABC$. Let $R\,$ and $S\,$ be the points where the circles inscribed in the triangles $ACH\,$ and $BCH^{}_{}$ are tangent to $\overline{CH}$. If $AB = 1995\,$, $AC = 1994\,$, and $BC = 1993\,$, then $RS\,$ can be expressed as $m/n\,$, where $m\,$ and $n\,$ are relatively prime integers. Find $m + n\,$.


 
 \textbf{Answer: }997



\textbf{Solution}


                    \begin{figure}[H]
                        \centering
                        \includegraphics[width=0.5\textwidth]{Solution/AIME_I_1993_15_Solution_Figure_1.png}
                    \end{figure}
                    

 From the \href{https://artofproblemsolving.com/wiki/index.php?title=Pythagorean\_Theorem}{Pythagorean Theorem}, $AH^2+CH^2=1994^2$, and $(1995-AH)^2+CH^2=1993^2$. 

 Subtracting those two equations yields $AH^2-(1995-AH)^2=3987$. 

 After simplification, we see that $2*1995AH-1995^2=3987$, or $AH=\frac{1995}{2}+\frac{3987}{2*1995}$.

 Note that $AH+BH=1995$. 

 Therefore we have that $BH=\frac{1995}{2}-\frac{3987}{2*1995}$. 

 Therefore $AH-BH=\frac{3987}{1995}$.

 
$\textbf{Note: }$An easier method is to use both facts together:

 First, by the Pythagorean Theorem, we find that $1993^2-BH^2=1994^2-AH^2$, so $AH^2-BH^2=1994^2-1993^2=3987$. We know that $AH+BH=1995$, so by difference of squares and dividing we find that $AH-BH=\frac{3987}{1995}=\frac{1329}{665}$. ~eevee9406

 
Now note that $RS=|HR-HS|$, $RH=\frac{AH+CH-AC}{2}$, and $HS=\frac{CH+BH-BC}{2}$. 

 Therefore we have $RS=\left| \frac{AH+CH-AC-CH-BH+BC}{2} \right|=\frac{|AH-BH-1994+1993|}{2}$.

 Plugging in $AH-BH$ and simplifying, we have $RS=\frac{1992}{1995*2}=\frac{332}{665} \rightarrow 332+665=\boxed{997}$.

 Edit by GameMaster402: 

 It can be shown that in any triangle with side lengths $n-1, n, n+1$, if you draw an altitude from the vertex to the side of $n+1$, and draw the incircles of the two right triangles, the distance between the two tangency points is simply $\frac{n-2}{2n+2}=\frac{n-2}{2(n+1)}$.

 Plugging in $n=1994$ yields that the answer is $\frac{1992}{2(1995)}$, which simplifies to $\frac{332}{665}$

 Edit by phoenixfire:

 It can further be shown for any triangle with sides $a=BC, b=CA, c=AB$ that
 \[RS=\dfrac{|b-a|}{2c}|a+b-c|\]Over here $a=1993, b=1994, c=1995$, so using the formula gives 
 \[RS = \dfrac{|1994 - 1993|}{2 \cdot 1995}|1993 + 1994 - 1995| = \dfrac{1 \cdot 1992}{2(1995)} = \frac{332}{665}.\]

 Note: We can also just write it as $RS=\frac{|b-a|(a+b-c)}{2c}$ since $a+b-c \geq 0$ by the triangle inequality. ~\href{https://artofproblemsolving.com/wiki/index.php?title=User:Yiyj1&action=edit&redlink=1}{Yiyj1}

 


\end{enumerate}